\documentclass[17pt]{article}

\usepackage{cmap}	% поиск по документу
\usepackage{hyperref}	% поддержка гиперссылок
\usepackage{geometry}	% настройка геометрии документа
\usepackage{multicol}   % настройка шаблона расположения текста
\usepackage{extsizes}	% расширенная настройка размеров шрифтов
\usepackage{indentfirst}    % отступ для первого абзаца
\usepackage{tabularx}	% расширенные таблицы
\usepackage{makecell}   % настройка клеток таблиц
\usepackage{caption}    % настройка подписей к вставкам
\usepackage[table]{xcolor}  % покраска таблиц
\usepackage{graphicx}	% вставка изображений
\usepackage{wrapfig, booktabs}	% оформление таблиц
\usepackage{graphicx, caption, subcaption}
\usepackage{minted}    % вставка кода

\usepackage{color}	% настройка цвета текста - \textcolor
\usepackage{soul}	% выделение текста - \hl
\usepackage[T2A]{fontenc}	% поддержка кириллических символов
\usepackage[utf8]{inputenc} % поддержка кириллических символов
\usepackage[english,russian]{babel}	% поддержка кириллических символов
\usepackage{csquotes, biblatex}   % вывод библиографии

\usepackage{mathcmd}	% математические команды
\usepackage{amsfonts}	% математические символы 1
\usepackage{amssymb}	% математические символы 2
\usepackage{amsmath}	% математические символы 3
\usepackage{amsthm}	% работа с теоремами

% Настройка окружения теорем 1 (RU)
\newtheorem{theorem}{Теорема}[section]
\newtheorem{corollary}{Следствие}[theorem]
\newtheorem{lemma}[theorem]{Лемма}
\newtheorem{example}{Пример}[section]
\newtheorem*{remark}{Утверждение}
\addto\captionsrussian{\renewcommand\proofname{Док-во:}}

% Настройка окружения теорем 2 (EN)
%\newtheorem{theorem}{TH}[section]
%\newtheorem{corollary}{CO}[theorem]
%\newtheorem{lemma}[theorem]{LM}
%\newtheorem{example}{EX}[section]
%\addto\captionsrussian{\renewcommand\proofname{PF:}}

\renewcommand{\labelitemii}{$\circ$}
\renewcommand{\labelitemiii}{-}

\addbibresource{biblio.bib}
\graphicspath{ {./images/} }

% Настройка оформления вставок кода
\setminted{xleftmargin=\parindent, fontsize=\footnotesize, linenos, tabsize=2, breaklines}
\usemintedstyle{emacs} 

% Настройка отступов
\geometry{
    a4paper,
    % total={200mm,287mm},
    left=10mm,
    right=10mm,
    top=10mm,
    bottom=20mm
}

% Заголовок
\title{
    \textbf{Методы моделирования и анализа стохастических систем} \\
    \large Лабораторная работа № 2: \\
    Нормальное распределение
}
\author{Костарев Иван (МЕНМ-150201)}
\date{весна 2025/2026}



\begin{document}

\maketitle
\tableofcontents
\newpage



\section{Метод 12}

Программная реализация метода 12:

\begin{minted}{python}
def random_normal_method_of_12(size: int = 1):
    n = 12
    uniform = rng.uniform(0., 1., size=(n, size)) - 0.5
    return np.sum(uniform, axis=0) * np.sqrt(12) / np.sqrt(n)    
\end{minted}

Смоделируем методом 12 нормально распределенную выборку $X \sim \mathcal{N}(0, 1)$ размера $N = 10^7$, и построим гистограмму плотности распределения (Рис. \ref{fig:method12_hist}).

\begin{figure}[htbp]
    \centering
    \includegraphics[width=0.6\textwidth]{method12_hist.png}
    \caption{Эмпирическая плотность распределения величины, полученной методом 12.}
    \label{fig:method12_hist}
\end{figure}

Программная реализация функции ошибки:

\begin{minted}{python}
def distribution_error(
    p: np.ndarray,
    p_true: np.ndarray
):
    return np.sum((p - p_true) ** 2)  
\end{minted}

Построим график зависимости погрешности \texttt{distribution\_error} от размера выборки $N$ (Рис. \ref{fig:method12_error}).

\begin{figure}[htbp]
    \centering
    \includegraphics[width=0.6\textwidth]{method12_error.png}
    \caption{Зависимость погрешности гистограммы от размера выборки $N$.}
    \label{fig:method12_error}
\end{figure}

Вычислим оценки мат ожидания и дисперсии для различных размеров выборки $N$, и сравним их с точными значениями для распределения $\mathcal{N}(0, 1)$ (Рис. \ref{fig:method12_means_plot}, \ref{fig:method12_vars_plot}).

\begin{figure}[htbp]
    \centering
    \includegraphics[width=0.6\textwidth]{method12_means_plot.png}
    \caption{Сравнение выборочного среднего с точным значением мат ожидания.}
    \label{fig:method12_means_plot}
\end{figure}

\begin{figure}[htbp]
    \centering
    \includegraphics[width=0.6\textwidth]{method12_vars_plot.png}
    \caption{Сравнение выборочной дисперсии с точным значением дисперсии.}
    \label{fig:method12_vars_plot}
\end{figure}

\textbf{Итог}: метод 12 для генерации нормально распределенной выборки требует довольно много памяти в текущей реализации на Python, поэтому эмпирические параметры распределения были вычислены для выборок размеров $N \leq 10^7$. Даже на таких размерах выборок графики демонстрируют, что с увеличением размера выборки мат ожидание и дисперсия стремятся к точным значениям. Ошибка плотности по сравнению с истинным нормальным распределением стремится к 0.



\newpage
\section{Преобразование Бокса -- Мюллера}

Программная реализация преобразования Бокса -- Мюллера:

\begin{minted}{python}
def random_normal_box_muller(size: int = 1):
    a, b = rng.uniform(0., 1., size=(2, size))
    xi_1 = np.sqrt(-2 * np.log(a)) * np.cos(2 * np.pi * b)
    xi_2 = np.sqrt(-2 * np.log(a)) * np.sin(2 * np.pi * b)
    return xi_1
\end{minted}

Смоделируем с помощью преобразования Бокса -- Мюллера нормально распределенную выборку $X \sim \mathcal{N}(0, 1)$ размера $N = 10^7$, и построим гистограмму плотности распределения (Рис. \ref{fig:box_muller_hist}).

\begin{figure}[htbp]
    \centering
    \includegraphics[width=0.6\textwidth]{box_muller_hist.png}
    \caption{Эмпирическая плотность распределения величины, полученной преобразованием Бокса -- Мюллера.}
    \label{fig:box_muller_hist}
\end{figure}

Построим график зависимости погрешности \texttt{distribution\_error} от размера выборки $N$ (Рис. \ref{fig:box_muller_error}).

\begin{figure}[htbp]
    \centering
    \includegraphics[width=0.6\textwidth]{box_muller_error.png}
    \caption{Зависимость погрешности гистограммы от размера выборки $N$.}
    \label{fig:box_muller_error}
\end{figure}

Вычислим оценки мат ожидания и дисперсии для различных размеров выборки $N$, и сравним их с точными значениями для распределения $\mathcal{N}(0, 1)$.

\begin{figure}[htbp]
    \centering
    \includegraphics[width=0.6\textwidth]{box_muller_means_plot.png}
    \caption{Сравнение выборочного среднего с точным значением мат ожидания.}
    \label{fig:box_muller_means_plot}
\end{figure}

\begin{figure}[htbp]
    \centering
    \includegraphics[width=0.6\textwidth]{box_muller_vars_plot.png}
    \caption{Сравнение выборочной дисперсии с точным значением дисперсии.}
    \label{fig:box_muller_vars_plot}
\end{figure}

\textbf{Итог}: преобразование Бокса -- Мюллера для генерации нормально распределенной выборки требует довольно много памяти в текущей реализации на Python, поэтому эмпирические параметры распределения были вычислены для выборок размеров $N \leq 10^7$. Даже на таких размерах выборок графики демонстрируют, что с увеличением размера выборки мат ожидание и дисперсия стремятся к точным значениям. Ошибка плотности по сравнению с истинным нормальным распределением стремится к 0.

\end{document}
